% Abhakseite „Das kann ich“ – Zone nur im Gesamt (\mitzone)
%% TODO Ich-kann-Sätze: Platzhalter wie in den Aufgabendateien
\begin{abhakseite}
\ifdefined\mitzone
\abhakgruppe{Kennst du schon}
\abhak{1}{Punkt vor Strich (3 · 4 − 5)}
\abhak{2}{Negative Zahlen addieren, subtrahieren, dividieren (12 : (−3))}
\abhak{3}{Terme zusammenfassen (gleichartige Glieder mit x, Zahl und x gemischt)}
\abhak{4}{Brüche: Hauptnenner, Bruch mal Zahl}
\abhak{5}{Klammern auflösen}
\abhak{6}{Terme zusammenfassen (gleichartige Glieder mit x, Zahl und x gemischt) – Fehler finden}
\abhak{7}{Terme zusammenfassen (gleichartige Glieder mit x, Zahl und x gemischt) – selbst rechnen}
\fi
\abhakgruppe{Einheit 1 · Gleichungen verstehen}
\abhak{8}{Umkehroperation benennen}
\abhak{9}{Lösung prüfen}
\abhak{10}{Durch Probieren lösen}
\abhak{11}{Lösung prüfen – Fehler finden, Begründen, Darstellungswechsel}
\abhakgruppe{Einheit 2 · Äquivalenzumformungen}
\abhak{12}{Vorzahl und Zeichen bei x lesen}
\abhak{13}{Reihenfolge bestimmen}
\abhak{14}{Umformen}
\abhak{15}{Umformen – Fehler finden, Begründen}
\abhakgruppe{Einheit 3 · Einheit 3}
\abhak{16}{x beidseitig}
\abhak{17}{Dezimalzahlen}
\abhak{18}{x beidseitig – Fehler finden, Begründen}
\abhakgruppe{Einheit 4 · Gleichungen aufstellen}
\abhak{19}{Aufstellen}
\abhak{20}{Formel umstellen}
\abhak{21}{Aufstellen – Fehler finden, Begründen, Darstellungswechsel, Anwendung}
\end{abhakseite}
